\section{Notation and Preliminaries}
\label{sec:prelim}

We study the allocation of $m$ indivisible items among $n$ agents, with
subsidies. Write $\agents = [n]$ for the set of agents and $\items$ for the set
of items, $\abs{\items} = m$. Each agent $i \in \agents$ has a valuation
$v_i : 2^{\items} \to \R$ with $v_i(\emptyset) = 0$, and we represent an
instance by the tuple $\langle \agents, \items, \set{v_i}_{i \in \agents}
\rangle$. Algorithms operate in the standard \emph{value-oracle} model: for each
$v_i$ we assume only an oracle that returns $v_i(S)$ for a queried set
$S \subseteq \items$, and running time is measured in $n$, $m$ and the number of
oracle calls.

An \emph{allocation} $\alloc = (\bundle{1}, \dots, \bundle{n})$ is an ordered
tuple of pairwise disjoint \emph{bundles}, with $\bundle{i}$ assigned to agent
$i$. The allocation is \emph{complete} if $\bigcup_{i \in \agents} \bundle{i} =
\items$ and \emph{partial} otherwise. The ordering matters throughout: much of
the theory below concerns permuting a fixed family of bundles among the agents,
and for a permutation $\sigma \in \mathbb{S}_n$ we write
$\alloc_\sigma := (A_{\sigma(1)}, \dots, A_{\sigma(n)})$ for the allocation in
which agent $i$ receives $A_{\sigma(i)}$. The utilitarian welfare of $\alloc$
is $\SW(\alloc) = \sum_{i \in \agents} v_i(\bundle{i})$.

\begin{definition}[Envy-Freeness]
\label{def:ef}
An allocation $\alloc = (\bundle{1},\dots,\bundle{n})$ is \emph{envy-free} (EF)
iff $v_i(\bundle{i}) \ge v_i(\bundle{j})$ for all agents $i, j \in \agents$.
\end{definition}

Envy-free allocations of indivisible items need not exist, so we allow a third
party to supplement the allocation with money. Agents have quasi-linear
utilities: agent $i$ receiving bundle $\bundle{i}$ and subsidy $\subsidy_i$
enjoys utility $v_i(\bundle{i}) + \subsidy_i$, with money entering linearly and
at the same exchange rate for every agent.

\begin{definition}[Envy-Free Solution]
\label{def:efsolution}
An allocation $\alloc = (\bundle{1},\dots,\bundle{n})$ and a subsidy vector
$\subsidy = (\subsidy_1,\dots,\subsidy_n) \in \R^n_{+}$ constitute an
\emph{envy-free solution} $(\alloc, \subsidy)$ iff
$v_i(\bundle{i}) + \subsidy_i \ge v_i(\bundle{j}) + \subsidy_j$ for all
$i, j \in \agents$.
\end{definition}

An allocation $\alloc$ is \emph{envy-freeable} if some $\subsidy \in \R^n_+$
makes $(\alloc,\subsidy)$ an envy-free solution; this is a property of the
allocation alone. Following~\cite{BKNS22} we write
$\Mset(\subsidy) := \set{i \in \agents \mid \subsidy_i \ge \subsidy_j \text{ for
all } j \in \agents}$ for the set of maximally subsidised agents.

We will also use the standard relaxation of envy-freeness up to one item (EF1). Since our framework allows chores, we use the two-sided form of EF1, in which one is allowed to remove at most one item from either of the two bundles being compared~\cite{KMSTY24}. Formally:
\begin{definition}[EF1, two-sided form]
\label{def:ef1}
An allocation $\alloc$ is \emph{envy-free up to one item} (EF1) iff for all
$i, j \in \agents$ there exists $X \subseteq \bundle{i} \cup \bundle{j}$ with
$\abs{X} \le 1$ such that
$v_i(\bundle{i} \setminus X) \ge v_i(\bundle{j} \setminus X)$.
\end{definition}

For an all-goods instance, removing an item from the envious agent's own bundle can only make the comparison harder, so the definition reduces to the usual goods version of EF1, in which one item may be removed from the envied bundle~\cite{HS19}. For an all-chore instance, the situation is reversed: removing a chore from the envied bundle can only make that bundle more attractive. Hence, if an item is removed, it may, without loss, be taken from the envious agent's own bundle.

Every instance in this paper is an all-chore instance, and there
Definition~\ref{def:ef1} takes a simpler form: the item removed can always
be taken from the envious agent's own bundle. We record that specialisation,
since it is the form that the results below establish.

\begin{definition}[EF1 for chores]
\label{def:ef1chores}
A complete allocation $\alloc$ of chores is \emph{EF1} iff for all
$i, j \in \agents$, either $i$ does not envy $j$, that is
$v_i(\bundle i) \ge v_i(\bundle j)$, or there is a chore
$e \in \bundle i$ with
\[
  v_i\bigl(\bundle i \setminus \set e\bigr) \;\ge\; v_i(\bundle j).
\]
\end{definition}

Definition~\ref{def:ef1chores} is exactly Definition~\ref{def:ef1} on an
all-chore instance. Taking $X = \emptyset$ in the first case and
$X = \set e \subseteq \bundle i$ in the second gives one direction; for the
converse, a removed item lying in $\bundle j$ only lowers
$\cost_i(\bundle j)$'s counterpart on the wrong side of the inequality, so
nothing is gained by removing from the envied bundle when every item is a
chore.

% --------------------------------------------------------------------------


\subsection{Negative Dichotomous Valuations}
\label{sec:negdich}

\begin{definition}[Negative dichotomous valuation] \label{def:negdich}
A valuation $v_i:2^{\items}\to\R$ with $v_i(\emptyset)=0$ is \emph{negative dichotomous} if, for every $S\subseteq\items$ and every $g\in\items\setminus S$,
\[
    \margplus{i}{S}{g}
    :=v_i(S\cup\set{g})-v_i(S)
    \in\set{-1,0}.
\]
Equivalently,
\[
    v_i(S)-v_i(S\cup\set{g})\in\set{0,1}.
\]
\end{definition}

Thus, adding an item weakly decreases an agent's valuation, so every item is a chore and $v_i$ is non-increasing:
\[
    S\subseteq T
    \quad\Longrightarrow\quad
    v_i(S)\ge v_i(T).
\]
Beyond the binary-marginal condition, we impose no further structure; in particular, the valuations need not be additive, submodular, or subadditive. This is the same level of generality considered by Barman et al.~\cite{BKNS22} for dichotomous goods.

For negative dichotomous valuations, it is often convenient to work with the corresponding non-negative \emph{cost function}
\[
    \cost_i(S):=-v_i(S).
\]
Then $\cost_i(\emptyset)=0$ and
\[
    \cost_i(S\cup\set{g})-\cost_i(S)\in\set{0,1}
\]
for every $S\subseteq\items$ and every $g\in\items\setminus S$.

\begin{definition}[Dichotomous cost function]
\label{def:dichcost}
A set function $\cost:2^{\items}\to\R$ is \emph{dichotomous} if $\cost(\emptyset)=0$ and
\[
    \cost(S\cup\set{g})-\cost(S)\in\set{0,1}
\]
for every $S\subseteq\items$ and every $g\in\items\setminus S$.
\end{definition}

Since all marginal costs are non-negative, every dichotomous cost function is automatically non-decreasing. Moreover, the correspondence
\[
    v_i \longleftrightarrow \cost_i=-v_i
\]
is a bijection between negative dichotomous valuations and dichotomous cost functions. In particular, the cost functions arising in our model belong to the same mathematical class as the dichotomous valuations studied by Barman et al.~\cite{BKNS22}; the difference is that here they represent costs rather than utilities.

In cost form, the envy-freeness condition for an allocation $\alloc=(\bundle1,\ldots,\bundle n)$ and subsidy vector $\subsidy$ becomes
\[
    \cost_i(\bundle i)-\subsidy_i
    \le
    \cost_i(\bundle j)-\subsidy_j
    \qquad\text{for all }i,j\in\agents.
\]
We use the valuation and cost representations interchangeably, according to which is more convenient.
% --------------------------------------------------------------------------
\subsection{The Envy Graph and the Halpern--Shah Characterisation}
\label{sec:envygraph}

We first fix the directed-graph terminology used below for the envy graph and,
in \Cref{sec:m-terminal}, for the equality graph of the partial-allocation
algorithm.

\begin{definition}[Strongly connected component; tail component]
\label{def:scc}
Let $G = (V,E)$ be a finite directed graph. Two vertices $u,v \in V$ are
\emph{strongly connected} if $G$ contains a directed path from $u$ to $v$ and a
directed path from $v$ to $u$. This relation is an equivalence relation on $V$,
and its equivalence classes are the \emph{strongly connected components} of $G$.
A strongly connected component $S$ is a \emph{tail component} (equivalently, a
\emph{sink component}) if no arc of $E$ leaves it, that is, if there is no arc
$(i,j) \in E$ with $i \in S$ and $j \in V \setminus S$.
\end{definition}

Being an equivalence relation, strong connectivity partitions $V$; hence every
vertex of $G$ lies in exactly one strongly connected component, and $G$ has at
least one such component whenever $V \ne \emptyset$. Tail components exist as
well.

\begin{observation}
\label{obs:tailscc}
Every finite directed graph $G = (V,E)$ with $V \ne \emptyset$ has at least one
tail strongly connected component.
\end{observation}

\begin{proof}
The \emph{condensation} of $G$ is the directed graph whose vertices are the
strongly connected components of $G$, with an arc from $S$ to $S'$ whenever
$S \ne S'$ and some arc of $E$ runs from a vertex of $S$ to a vertex of $S'$.
The condensation is acyclic: a directed cycle through two distinct components
would place every vertex on it in a single component, contradicting their
distinctness~\cite[Ch.~1]{BJG09}. Since $V \ne \emptyset$, the condensation has
at least one vertex, and a finite acyclic directed graph on a non-empty vertex
set has a vertex of out-degree $0$ --- follow arcs from any vertex, never
repeating one, and the walk must terminate. A component with out-degree $0$ in
the condensation has no arc of $E$ leaving it, so it is a tail component.
\end{proof}


For an allocation $\alloc$, the \emph{envy graph} $\envygraph{\alloc}$ is the
complete weighted directed graph on vertex set $\agents$ in which the arc
$(i,j)$ carries weight
\[
  \edgew{\alloc}(i,j) \;:=\; v_i(\bundle{j}) - v_i(\bundle{i})
  \;=\; \cost_i(\bundle{i}) - \cost_i(\bundle{j}),
\]
the envy that agent $i$ has towards agent $j$. The weight of a directed path
$P$ is $\edgew{\alloc}(P) = \sum_{(i,j) \in P} \edgew{\alloc}(i,j)$, and
$\pathw{\alloc}(i)$ denotes the maximum weight of any path in
$\envygraph{\alloc}$ starting at $i$, the empty path of weight $0$ included.

The following two results of Halpern and Shah~\cite{HS19} underpin much of the literature on subsidies for envy-freeness. They apply to general valuation functions and hence, in particular, to the negative dichotomous valuations considered here. We state them in the notation of our model.
\begin{theorem}[\cite{HS19}]
\label{thm:hs-characterisation}
For any allocation $\alloc = (\bundle{1},\dots,\bundle{n})$, the following are
equivalent.
\begin{enumerate}
  \item[(i)] $\alloc$ is envy-freeable.
  \item[(ii)] $\alloc$ maximises utilitarian welfare across all reassignments of
    its own bundles among the agents: for every permutation $\sigma$ over
    $\agents$, $\sum_{i \in \agents} v_i(\bundle{i}) \ge
    \sum_{i \in \agents} v_i(A_{\sigma(i)})$.
  \item[(iii)] The envy graph $\envygraph{\alloc}$ has no positive-weight
    directed cycle.
\end{enumerate}
\end{theorem}

\begin{theorem}[\cite{HS19}]
\label{thm:hs-minsubsidy}
For any envy-freeable allocation $\alloc$, the subsidy vector $\optsubsidy$
given by $\optsubsidy_i := \pathw{\alloc}(i)$ realises an envy-free solution
$(\alloc, \optsubsidy)$, and $\optsubsidy_i \le \subsidy_i$ for every
$i \in \agents$ and every $\subsidy$ such that $(\alloc,\subsidy)$ is an
envy-free solution. It can be computed in strongly polynomial time by running
an all-pairs longest-path computation on $\envygraph{\alloc}$.
\end{theorem}

Condition~(ii) of Theorem~\ref{thm:hs-characterisation} implies that, starting from any collection of bundles, one can obtain an envy-freeable allocation by suitably reassigning those bundles among the agents. In particular, such a reassignment can be found by solving a maximum-weight perfect matching problem between agents and bundles. Accordingly, the main issue is not the existence of an envy-freeable allocation, but rather the magnitude of the subsidies required to support one, as characterised by Theorem~\ref{thm:hs-minsubsidy}. 


For negative dichotomous valuations, all bundle values are integral:
indeed, $v_i(S)$ is obtained by telescoping from $v_i(\emptyset)=0$
along marginals in $\{-1,0\}$. Consequently, every arc weight
$\edgew{\alloc}(i,j)$ in the envy graph is an integer, and hence, by
Theorem~\ref{thm:hs-minsubsidy}, the pointwise-minimal subsidy vector
$\optsubsidy$ is integral as well. Therefore, showing
$\optsubsidy_i\leq 1$ for every $i$ is equivalent to showing
$\optsubsidy\in\{0,1\}^n$.

Moreover, every pointwise-minimal nonnegative subsidy vector has at
least one zero coordinate. Indeed, if $\optsubsidy_i>0$ for every
$i$, then subtracting $\min_i\optsubsidy_i$ from every coordinate
preserves all envy-freeness inequalities while producing a strictly
smaller nonnegative subsidy vector, contradicting pointwise minimality.
Thus, if $\optsubsidy\in\{0,1\}^n$, then automatically
\[
    \sum_{i\in\agents}\optsubsidy_i\leq n-1.
\]
% Accordingly, for the pointwise-minimal subsidy vector, our target can
% be written simply as
% \[
%     \exists\,\alloc:\qquad
%     \max_{i\in\agents}\pathw{\alloc}(i)\leq 1.
% \]





% \begin{observation}[Integrality]
% \label{obs:integrality}
% Let $\set{v_i}_{i \in \agents}$ be negative dichotomous. Then $v_i(S) \in
% \Z_{\le 0}$ for every $i$ and every $S \subseteq \items$; every arc weight
% $\edgew{\alloc}(i,j)$ is an integer; and for every envy-freeable allocation
% $\alloc$ the minimum subsidy vector satisfies $\optsubsidy \in \Z^n_{\ge 0}$.
% \end{observation}

% \begin{proof}
% Fix $i$ and $S = \set{g_1,\dots,g_k}$. Telescoping from $v_i(\emptyset) = 0$
% along any ordering of $S$ writes $v_i(S)$ as a sum of $k$ marginals, each of
% which lies in $\set{-1,0}$ by Definition~\ref{def:negdich}; hence $v_i(S) \in \Z$ and
% $-k \le v_i(S) \le 0$. Consequently each arc weight
% $\edgew{\alloc}(i,j) = v_i(\bundle{j}) - v_i(\bundle{i})$ is a difference of
% integers, and the weight of any directed path, being a finite sum of arc
% weights, is an integer. By Theorem~\ref{thm:hs-minsubsidy},
% $\optsubsidy_i = \pathw{\alloc}(i)$ is the maximum such weight over paths
% starting at $i$, and it is non-negative because the empty path has weight $0$.
% \end{proof}

% Observation~\ref{obs:integrality} is what makes the target statement --- a subsidy vector in
% $\set{0,1}^n$ --- a statement about a \emph{bound} rather than about rounding:
% once $\optsubsidy$ is known to be integral, showing $\optsubsidy_i \le 1$ for
% every $i$ is the same as showing $\optsubsidy \in \set{0,1}^n$. The same
% observation underlies the corresponding step in~\cite{BKNS22}, where the arc weights of a dichotomous instance are likewise integers.

% The following observation further simplifies the target guarantee.

% \begin{proposition}[Some agent is always unpaid]
% \label{prop:zero-coordinate}
% For every envy-freeable allocation $\alloc$ there is an agent $h$ with
% $\optsubsidy_h = 0$. Consequently, if $\optsubsidy \in \set{0,1}^n$ then
% $\sum_{i \in \agents} \optsubsidy_i \le n-1$.
% \end{proposition}

% \begin{proof}
% By Theorem~\ref{thm:hs-characterisation}(iii) the envy graph $\envygraph{\alloc}$
% has no positive-weight cycle, so a walk of maximum weight may be taken simple
% and $\pathw{\alloc}$ is finite; it is non-negative because the empty path
% qualifies. Choose a path $P$ attaining $\max_{i} \pathw{\alloc}(i)$ and let $h$
% be its final vertex. If $\pathw{\alloc}(h) > 0$ there is a positive-weight walk
% $Q$ leaving $h$, and $P$ followed by $Q$ is a walk from $P$'s start of weight
% $\edgew{\alloc}(P) + \edgew{\alloc}(Q) > \edgew{\alloc}(P)$, contradicting the
% maximality of $P$ among walks. Hence $\pathw{\alloc}(h) = 0$. The second claim is immediate: at
% most $n-1$ coordinates of $\optsubsidy$ can equal $1$.
% \end{proof}

% Proposition~\ref{prop:zero-coordinate} justifies the word ``hence'' in the
% target statement of \Cref{sec:intro}: the bound on the \emph{total} is not an
% independent requirement but a consequence of the per-agent one. For the
% pointwise-minimal subsidy the target is therefore the single-clause statement
% \[
%   \exists\, \alloc \ :\quad \max_{i \in \agents} \pathw{\alloc}(i) \;\le\; 1 .
% \]


% Two cautions are worth recording alongside it. First, the bound cannot be improved: the family of
% Example~\ref{ex:lowerbound} attains $n-1$ for every $n$, verified exhaustively
% for $n \le 7$. Second, the conjecture is \emph{not} the assertion that the
% utilitarian-optimal subsidy is at most $n-1$. Minimising the total over all
% allocations can do strictly better than any allocation with
% $\optsubsidy \in \set{0,1}^n$, by concentrating the subsidy on a single agent:
% there are instances at $n = 4$ whose cheapest envy-free solution costs $2$, borne
% entirely by one agent, while every allocation with $\optsubsidy \in \set{0,1}^4$
% costs $3$. The $\set{0,1}$ shape is a genuine restriction, not a normalisation of
% the total.

